% ======================================================================
% 3. FRAMING  (budget: 1.5 pages)
% Job of the chapter: give the reader one mental model that makes PtH, PtK
% (and the two out-of-scope variants) the same shape, then pin down exactly
% what an attacker must already own.
% ======================================================================
\section{Credential Material as a Bearer Capability}
\label{sec:framing}

\subsection{The unifying property}
% CLAIM (0.4): the authenticator is any secret from which the protocol's
% proof can be computed -- not the password.
% Beats: enumerate what Windows may hold (NT hash, AES keys, ticket session
% keys, PKINIT certificates); each material has a matching "pass-the-X";
% all of them die to the same class of defence (material must not exist on
% the host, or must be unusable off it).

\subsection{Threat model and preconditions}
\subsubsection{Attacker starting positions}
% CLAIM (0.3): two entry points only -- local admin on one host, or material
% obtained out of band. Everything else in the paper follows from these.
\subsubsection{The four preconditions}
% CLAIM (0.4): hash + a service that accepts NTLM + authorisation on the
% target + reachability. Removing any one of them ends the attack, which is
% what chapter 9 exploits.
\subsubsection{Lateral movement, not privilege escalation}
% CLAIM (0.3): PtH transports an existing privilege set; it does not create
% one. Include the local-account caveat: UAC remote restrictions and
% LocalAccountTokenFilterPolicy, the most common reason a real-world attempt
% fails at the last step.

\subsection{Taxonomy and position of this study}
% CLAIM (0.3): a small table positioning PtH, PtK, and the two excluded
% variants by material consumed and protocol family spoken. One paragraph on
% why the excluded two need a different lab (ticket lifecycle, AD CS).
